\section{Билет 28. Энергия упругой волны с выводом. Плотность энергии. Среднее значение плотности энергии.}

\begin{definition}
Волновой процесс сопровождается переносом энергии. Для количественной характеристики вводится \textbf{объемная плотность энергии волны} $w$ --- энергия, приходящаяся на единицу объема среды:
\begin{equation}
    w = \frac{dW}{dV} = w_k + w_p,
\end{equation}
где $w_k$ --- плотность кинетической энергии, $w_p$ --- плотность потенциальной энергии упругой деформации.
\end{definition}

\begin{theorem}[Вывод плотности энергии плоской упругой волны]
Рассмотрим плоскую монохроматическую волну, распространяющуюся вдоль оси $x$:
\begin{equation}
    \xi(x,t) = A \cos(\omega t - kx). \label{eq:wave_28}
\end{equation}
Выделим в среде элементарный объем $dV = S \, dx$ с массой $dm = \rho \, dV$, где $\rho$ --- плотность среды.

\textbf{1. Кинетическая энергия.} Скорость колебаний частиц в данном объеме равна производной смещения по времени:
\begin{equation}
    v = \frac{\partial \xi}{\partial t} = -A\omega \sin(\omega t - kx).
\end{equation}
Кинетическая энергия элемента:
\begin{equation}
    dW_k = \frac{1}{2} dm \, v^2 = \frac{1}{2} \rho S \, dx \, A^2 \omega^2 \sin^2(\omega t - kx).
\end{equation}
Плотность кинетической энергии:
\begin{equation}
    w_k = \frac{dW_k}{dV} = \frac{1}{2} \rho A^2 \omega^2 \sin^2(\omega t - kx). \label{eq:wk}
\end{equation}

\textbf{2. Потенциальная энергия.} Для упругой среды потенциальная энергия определяется работой сил упругости при деформации. Относительная деформация элемента $\varepsilon = \frac{\partial \xi}{\partial x} = Ak \sin(\omega t - kx)$. Согласно закону Гука, плотность потенциальной энергии равна:
\begin{equation}
    w_p = \frac{1}{2} E \varepsilon^2 = \frac{1}{2} E A^2 k^2 \sin^2(\omega t - kx),
\end{equation}
где $E$ --- модуль упругости среды. Используя связь $v = \sqrt{E/\rho}$ и $k = \omega/v$, получаем $E k^2 = \rho v^2 (\omega/v)^2 = \rho \omega^2$. Следовательно:
\begin{equation}
    w_p = \frac{1}{2} \rho A^2 \omega^2 \sin^2(\omega t - kx). \label{eq:wp}
\end{equation}

\textbf{3. Полная плотность энергии.} Суммируя \eqref{eq:wk} и \eqref{eq:wp}, находим:
\begin{equation}
    w(x,t) = w_k + w_p = \rho A^2 \omega^2 \sin^2(\omega t - kx). \label{eq:total_w}
\end{equation}
\end{theorem}

\begin{property}
Из формулы \eqref{eq:total_w} следуют важные выводы:
\begin{enumerate}
    \item В бегущей волне плотности кинетической и потенциальной энергий равны в каждой точке и в каждый момент времени: $w_k = w_p$.
    \item Полная плотность энергии не постоянна, а пульсирует с частотой $2\omega$, достигая максимума $w_{\max} = \rho A^2 \omega^2$ в моменты прохождения частицами положения равновесия и обращаясь в ноль в точках максимального отклонения.
    \item Энергия переносится волной со скоростью $v$.
\end{enumerate}
\end{property}

\begin{definition}[Среднее значение плотности энергии]
Поскольку плотность энергии быстро изменяется со временем, на практике используют её среднее значение за период колебаний $T$. Учитывая, что среднее значение $\sin^2$ за период равно $1/2$:
\begin{equation}
    \langle \sin^2(\omega t - kx) \rangle = \frac{1}{T} \int_0^T \sin^2(\omega t - kx) \, dt = \frac{1}{2},
\end{equation}
получаем:
\begin{equation}
    \langle w \rangle = \frac{1}{2} \rho A^2 \omega^2. \label{eq:avg_w}
\end{equation}
Средняя плотность энергии пропорциональна квадрату амплитуды и квадрату частоты волны. Эта величина является ключевой для определения интенсивности волны и вектора Умова (плотности потока энергии).
\end{definition}

\begin{remark}
Для стоячей волны распределение энергии имеет иной характер: энергия не переносится вдоль среды, а периодически переходит из кинетической формы в потенциальную и обратно внутри каждого сегмента между узлами. В моменты максимального смещения вся энергия сосредоточена в потенциальной форме (максимальная деформация у узлов, нулевая у пучностей), а при прохождении через равновесие --- в кинетической (максимальная скорость у пучностей). Среднее по времени значение плотности энергии стоячей волны в пространстве также равно $\frac{1}{2} \rho A^2 \omega^2$, но локально оно распределено неоднородно.
\end{remark}
